% Source: https://artofproblemsolving.com/community/c3h2500092_are_negatives_real
% Stable identity: aops-topic-2500092-q01
\textbf{(a)} If $x$ is a whole number, then $x\ge0$, so $b+x\ge b>a$ and $b+x\ne a$: there is no solution among the whole numbers. Among the integers, adding $-b$ to both sides shows that $b+x=a$ holds exactly when $x=a-b$. So there is exactly one integer solution, $x=a-b$, and it is negative because $a<b$. For example, $5+x=3$ has the unique solution $x=3-5=\boxed{-2}$.

\textbf{(b)} The argument shows only that a negative number cannot be the answer to the question ``how many objects are there?'' But numbers are used for more than counting collections. Negative numbers describe quantities measured from a reference point or with a direction: a temperature $2$ degrees below zero is $-2$ degrees, a debt of $2$ dollars is a balance of $-2$ dollars, a fall of $2$ metres is a change in height of $-2$ metres, and the point $2$ units to the left of $0$ on the number line is $-2$. In mathematics, $-1$ is the number that gives $0$ when it is added to $1$, and part (a) shows why such numbers are needed: without them, even an equation as simple as $5+x=3$ has no solution. The real numbers are the numbers of the whole number line, negative ones included; the word ``real'' is the name of this number system, not a claim that every number counts physical objects.
